% =====================================================================
\section{Relation to prior work}
\label{ch:related}
% =====================================================================

This chapter positions the campaign relative to the main bodies of work on the Yang--Mills mass
gap. It is not a survey; it is a map of what the paper imports, what it does not claim, and where
its contribution sits. The corpus's own surveys (\code{SURVEY-ADJACENT-TOOLS.md},
\code{SURVEY-CONSTRUCTIVE.md}, \code{SURVEY-RECENT-ARXIV.md}, \code{SURVEY-AI-FORMAL.md},
\code{SURVEY-WILDCARD-DOMAINS.md}, \code{SURVEY-MASTER-SYNTHESIS.md},
\code{docs/yang-mills-prior-work.md}) contain the detailed reading notes; this chapter states the
conclusions.

\subsection{Constructive field theory}

The constructive programme asks for a measure on the space of gauge fields satisfying the
Osterwalder--Schrader axioms. Its main achievements are ultraviolet: Balaban's renormalization
group analysis of lattice gauge theory \cite{balaban}, and the infrared construction with a
cutoff by Magnen--Rivasseau--S\'en\'eor \cite{mrs}. The modern probabilistic approach
(Bauerschmidt--Brydges--Slade for $\varphi^4_3$ \cite{bbs}; Chatterjee's rigorous large-$N$ and
survey \cite{chatterjee}) and the stochastic-quantization approach (regularity structures) are
powerful but remain small-parameter theories in the relevant sense. None of them produces a
uniform mass gap at the physical four-dimensional coupling; the obstruction is the same
intermediate-coupling region that this paper identifies as $(G1)$/H1. The paper's G1 node board
(Chapter~\ref{ch:existence}) is organized around exactly this split: the perturbative kernel is
closed, the ultraviolet node L3 is open, and the infrared node is the gap.

\subsection{Lattice gauge theory}

The lattice is where the mass gap is measured. Wilson's formulation \cite{wilson} and Creutz's
Monte Carlo study \cite{creutz} established the framework; the strong-coupling expansion and its
rigorous polymer form (Osterwalder--Seiler \cite{osterwalder-seiler}, the source of P3's
technology) give the only unconditional gap theorems, at $\beta<1/48$. Shen--Zhu--Zhu
\cite{szz} extended the rigorous strong-coupling gap to a modern standard. On the numerical side,
the glueball spectrum (Morningstar--Peardon \cite{morningstar}), the step-scaling method
\cite{luscher-stepscaling}, the gradient flow \cite{luscher-flow}, and the effective string
description of the flux tube \cite{luscher-weisz, polchinski-strominger} are the standard tools.
This paper's numerics (Chapter~\ref{ch:numerics}) are small by comparison, and are presented as
diagnostics rather than as a spectrum computation. The campaign's contribution on this side is
the diagnostic set (P3 approach law, Re-slope, finite-size law) and the explicit statement of
what a decisive computation would require.

\subsection{Holography and AdS/QCD}

The holographic approach constructs confining gauge theories from string backgrounds
(Maldacena \cite{maldacena}; Gubser--Klebanov--Polyakov \cite{gkp}; Witten \cite{witten,
witten-thermal}), with hard-wall and soft-wall models (Polchinski--Strassler \cite{polchinski-strassler};
Sakai--Sugimoto \cite{sakai-sugimoto}) reproducing the qualitative glueball spectrum. The
campaign's holographic wave (Chapter~\ref{ch:betarg}, the NEC route) found the precise limitation:
the bulk null-energy condition is compatible with a gapless infrared AdS throat, and a confining
wall requires a modeling choice whose justification is the dynamical question. Holography gives
a dictionary and a picture, not a proof for the four-dimensional theory.

\subsection{The conformal bootstrap and CFT data}

The conformal bootstrap (El-Showk et al.\ \cite{elshowk}; the review of Poland--Rychkov--Vichi
\cite{prv}) has transformed the study of CFTs by turning crossing symmetry and unitarity into
numerical bounds. Its relation to the mass gap is the subject of Chapter~\ref{ch:betarg}: the
bootstrap tests the \emph{existence} of a CFT with given operator content, whereas the gap is a
statement about the \emph{flow}. The conformal-collider bounds on $a/c$
\cite{hofman-maldacena, hofman-li} and the $a$-theorem \cite{komargodski-schwimmer} are exact constraints that are
compatible with a gapless infrared; they constrain which CFT, not whether one is reached. The
campaign's conclusion is that the bootstrap is necessary but not sufficient, and that a pure-glue
bootstrap, if it existed, would exclude one candidate endpoint rather than prove the gap.

\subsection{Resurgence and transseries}

The resurgent transseries programme ('t Hooft's renormalons; the review of
Aniceto--Başar--Schiappa \cite{abs}) is the natural language for the essential singularity
$e^{-1/(2\beta_0g^2)}$. The campaign's analytic-transfer no-go (Chapter~\ref{sec:esssing}) is a
statement in this language: the nonperturbative scale is invisible to every analytic-in-$g^2$
procedure, and the condensate's $u=2$ renormalon ambiguity is the same size as the condensate.
Resurgence supplies the correct bookkeeping; it does not supply the sign.

\subsection{Large-$N$ and loop equations}

The large-$N$ limit ('t Hooft \cite{thooft-planar}; Makeenko--Migdal \cite{makeenko-migdal};
Eguchi--Kawai \cite{eguchi-kawai}) organizes the theory around a master field and loop equations.
Chatterjee's rigorous results \cite{chatterjee} make part of this picture a theorem at strong
coupling. The campaign's large-$N$ route (Chapter~\ref{ch:betarg}) found that $1/N$ controls
fluctuations, not the fixed-point question: the 't Hooft coupling is $O(1)$ at the confining
scale, and the four-dimensional master field is not constructed.

\subsection{Spectral, geometric, and dependence methods}

The campaign also drew on spectral theory (Kato \cite{kato}; Reed--Simon \cite{reed-simon}),
Brascamp--Lieb and Pr\'ekopa \cite{bl,prekopa}, the theory of negative dependence
(FKG \cite{fkg}; Borcea--Br\"and\'en--Liggett \cite{bbl}), optimal transport
(von Renesse--Sturm \cite{vrs}), and the Perron--Frobenius/Krein--Rutman theory. These are the
tools of Part~I. Their collective lesson is the Schur-convex master obstruction
(Chapter~\ref{sec:schur-master}): the standard positivity and convexity toolkit is blind to the
atom, because the free field is a feasible atomless point. The campaign's no-gos are statements
about these classes, not about particular papers.

\subsection{Formalization and AI-assisted mathematics}

The Lean~4 formalization (Appendix~\ref{app:lean}) uses the pinned toolchain and Mathlib. The
campaign's contribution is not a new formalization technology but a disciplined use of one: the
reduction's logic is machine-checked, the H1 sharpening is machine-checked, and the analytic
hypotheses are explicitly displayed as \code{sorry} rather than hidden. The AI contribution
statement (Chapter~\ref{sec:verdict}) records what the AI systems did and did not do. The
paper's position is that formalization is most valuable for the \emph{logical skeleton} --- the
implication structure and the countermodels --- and least valuable for the analytic content, which
remains a matter of proof and of labelled conjecture.

\subsection{What this paper adds, and what it does not}

\textbf{Adds.} (i) A labelled inventory of what can and cannot be concluded, with a
machine-checked core. (ii) The two-coordinate picture: the wall in the measure/weak-$*$
coordinate and in the $\beta$/RG/conformal coordinate, with disjoint technique pools and
independent diagnoses. (iii) The weak-$*$ theorem and its consequences for the class of forcing
functionals. (iv) The $(\star)$ audit and the $(G1)$ classification, including the reality of the
conformal window and the insufficiency of each modern technique. (v) The Savvidy cleanup: a false
obstruction removed by theorems, with the surviving margin precisely located. (vi) The diagnostic
set for the atom. (vii) An errata record, including a corrected central-charge normalization.

\textbf{Does not add.} (i) A proof of the gap. (ii) A new constructive field theory. (iii) A new
lattice computation at the precision needed for a one-sided bound. (iv) A solution of $(G1)$ or of
the negative-association inequality. The paper is an audit and a specification, not a
breakthrough, and it says so on every page.
